\batchmode
\documentclass[12pt]{article}
\RequirePackage{ifthen}


\usepackage{amsmath}
\usepackage{amssymb}
\usepackage[french]{babel}
\usepackage[latin1]{inputenc}
\usepackage[T1]{fontenc}
\usepackage[all]{xy}
\usepackage{html}

%
\providecommand{\classe}[1]{\,\begin{tabular}{|l|}
\!\!  #1  \!\!\! \\
\hline
\end{tabular}\:} 

%
\providecommand{\der}[1]{\overset{\textrm{\footnotesize (#1)\;\;\;}}{\Longrightarrow}} 




\usepackage[dvips]{color}


\pagecolor[gray]{.7}

\usepackage[latin1]{inputenc}



\makeatletter

\makeatletter
\count@=\the\catcode`\_ \catcode`\_=8 
\newenvironment{tex2html_wrap}{}{}%
\catcode`\<=12\catcode`\_=\count@
\newcommand{\providedcommand}[1]{\expandafter\providecommand\csname #1\endcsname}%
\newcommand{\renewedcommand}[1]{\expandafter\providecommand\csname #1\endcsname{}%
  \expandafter\renewcommand\csname #1\endcsname}%
\newcommand{\newedenvironment}[1]{\newenvironment{#1}{}{}\renewenvironment{#1}}%
\let\newedcommand\renewedcommand
\let\renewedenvironment\newedenvironment
\makeatother
\let\mathon=$
\let\mathoff=$
\ifx\AtBeginDocument\undefined \newcommand{\AtBeginDocument}[1]{}\fi
\newbox\sizebox
\setlength{\hoffset}{0pt}\setlength{\voffset}{0pt}
\addtolength{\textheight}{\footskip}\setlength{\footskip}{0pt}
\addtolength{\textheight}{\topmargin}\setlength{\topmargin}{0pt}
\addtolength{\textheight}{\headheight}\setlength{\headheight}{0pt}
\addtolength{\textheight}{\headsep}\setlength{\headsep}{0pt}
\setlength{\textwidth}{349pt}
\newwrite\lthtmlwrite
\makeatletter
\let\realnormalsize=\normalsize
\global\topskip=2sp
\def\preveqno{}\let\real@float=\@float \let\realend@float=\end@float
\def\@float{\let\@savefreelist\@freelist\real@float}
\def\liih@math{\ifmmode$\else\bad@math\fi}
\def\end@float{\realend@float\global\let\@freelist\@savefreelist}
\let\real@dbflt=\@dbflt \let\end@dblfloat=\end@float
\let\@largefloatcheck=\relax
\let\if@boxedmulticols=\iftrue
\def\@dbflt{\let\@savefreelist\@freelist\real@dbflt}
\def\adjustnormalsize{\def\normalsize{\mathsurround=0pt \realnormalsize
 \parindent=0pt\abovedisplayskip=0pt\belowdisplayskip=0pt}%
 \def\phantompar{\csname par\endcsname}\normalsize}%
\def\lthtmltypeout#1{{\let\protect\string \immediate\write\lthtmlwrite{#1}}}%
\newcommand\lthtmlhboxmathA{\adjustnormalsize\setbox\sizebox=\hbox\bgroup\kern.05em }%
\newcommand\lthtmlhboxmathB{\adjustnormalsize\setbox\sizebox=\hbox to\hsize\bgroup\hfill }%
\newcommand\lthtmlvboxmathA{\adjustnormalsize\setbox\sizebox=\vbox\bgroup %
 \let\ifinner=\iffalse \let\)\liih@math }%
\newcommand\lthtmlboxmathZ{\@next\next\@currlist{}{\def\next{\voidb@x}}%
 \expandafter\box\next\egroup}%
\newcommand\lthtmlmathtype[1]{\gdef\lthtmlmathenv{#1}}%
\newcommand\lthtmllogmath{\dimen0\ht\sizebox \advance\dimen0\dp\sizebox
  \ifdim\dimen0>.95\vsize
   \lthtmltypeout{%
*** image for \lthtmlmathenv\space is too tall at \the\dimen0, reducing to .95 vsize ***}%
   \ht\sizebox.95\vsize \dp\sizebox\z@ \fi
  \lthtmltypeout{l2hSize %
:\lthtmlmathenv:\the\ht\sizebox::\the\dp\sizebox::\the\wd\sizebox.\preveqno}}%
\newcommand\lthtmlfigureA[1]{\let\@savefreelist\@freelist
       \lthtmlmathtype{#1}\lthtmlvboxmathA}%
\newcommand\lthtmlpictureA{\bgroup\catcode`\_=8 \lthtmlpictureB}%
\newcommand\lthtmlpictureB[1]{\lthtmlmathtype{#1}\egroup
       \let\@savefreelist\@freelist \lthtmlhboxmathB}%
\newcommand\lthtmlpictureZ[1]{\hfill\lthtmlfigureZ}%
\newcommand\lthtmlfigureZ{\lthtmlboxmathZ\lthtmllogmath\copy\sizebox
       \global\let\@freelist\@savefreelist}%
\newcommand\lthtmldisplayA{\bgroup\catcode`\_=8 \lthtmldisplayAi}%
\newcommand\lthtmldisplayAi[1]{\lthtmlmathtype{#1}\egroup\lthtmlvboxmathA}%
\newcommand\lthtmldisplayB[1]{\edef\preveqno{(\theequation)}%
  \lthtmldisplayA{#1}\let\@eqnnum\relax}%
\newcommand\lthtmldisplayZ{\lthtmlboxmathZ\lthtmllogmath\lthtmlsetmath}%
\newcommand\lthtmlinlinemathA{\bgroup\catcode`\_=8 \lthtmlinlinemathB}
\newcommand\lthtmlinlinemathB[1]{\lthtmlmathtype{#1}\egroup\lthtmlhboxmathA
  \vrule height1.5ex width0pt }%
\newcommand\lthtmlinlineA{\bgroup\catcode`\_=8 \lthtmlinlineB}%
\newcommand\lthtmlinlineB[1]{\lthtmlmathtype{#1}\egroup\lthtmlhboxmathA}%
\newcommand\lthtmlinlineZ{\egroup\expandafter\ifdim\dp\sizebox>0pt %
  \expandafter\centerinlinemath\fi\lthtmllogmath\lthtmlsetinline}
\newcommand\lthtmlinlinemathZ{\egroup\expandafter\ifdim\dp\sizebox>0pt %
  \expandafter\centerinlinemath\fi\lthtmllogmath\lthtmlsetmath}
\newcommand\lthtmlindisplaymathZ{\egroup %
  \centerinlinemath\lthtmllogmath\lthtmlsetmath}
\def\lthtmlsetinline{\hbox{\vrule width.1em \vtop{\vbox{%
  \kern.1em\copy\sizebox}\ifdim\dp\sizebox>0pt\kern.1em\else\kern.3pt\fi
  \ifdim\hsize>\wd\sizebox \hrule depth1pt\fi}}}
\def\lthtmlsetmath{\hbox{\vrule width.1em\kern-.05em\vtop{\vbox{%
  \kern.1em\kern0.8 pt\hbox{\hglue.17em\copy\sizebox\hglue0.8 pt}}\kern.3pt%
  \ifdim\dp\sizebox>0pt\kern.1em\fi \kern0.8 pt%
  \ifdim\hsize>\wd\sizebox \hrule depth1pt\fi}}}
\def\centerinlinemath{%
  \dimen1=\ifdim\ht\sizebox<\dp\sizebox \dp\sizebox\else\ht\sizebox\fi
  \advance\dimen1by.5pt \vrule width0pt height\dimen1 depth\dimen1 
 \dp\sizebox=\dimen1\ht\sizebox=\dimen1\relax}

\def\lthtmlcheckvsize{\ifdim\ht\sizebox<\vsize 
  \ifdim\wd\sizebox<\hsize\expandafter\hfill\fi \expandafter\vfill
  \else\expandafter\vss\fi}%
\providecommand{\selectlanguage}[1]{}%
\makeatletter \tracingstats = 1 
\providecommand{\Beta}{\textrm{B}}
\providecommand{\Mu}{\textrm{M}}
\providecommand{\Kappa}{\textrm{K}}
\providecommand{\Rho}{\textrm{R}}
\providecommand{\Epsilon}{\textrm{E}}
\providecommand{\Chi}{\textrm{X}}
\providecommand{\Iota}{\textrm{J}}
\providecommand{\omicron}{\textrm{o}}
\providecommand{\Zeta}{\textrm{Z}}
\providecommand{\Eta}{\textrm{H}}
\providecommand{\Nu}{\textrm{N}}
\providecommand{\Omicron}{\textrm{O}}
\providecommand{\Tau}{\textrm{T}}
\providecommand{\Alpha}{\textrm{A}}


\begin{document}
\pagestyle{empty}\thispagestyle{empty}\lthtmltypeout{}%
\lthtmltypeout{latex2htmlLength hsize=\the\hsize}\lthtmltypeout{}%
\lthtmltypeout{latex2htmlLength vsize=\the\vsize}\lthtmltypeout{}%
\lthtmltypeout{latex2htmlLength hoffset=\the\hoffset}\lthtmltypeout{}%
\lthtmltypeout{latex2htmlLength voffset=\the\voffset}\lthtmltypeout{}%
\lthtmltypeout{latex2htmlLength topmargin=\the\topmargin}\lthtmltypeout{}%
\lthtmltypeout{latex2htmlLength topskip=\the\topskip}\lthtmltypeout{}%
\lthtmltypeout{latex2htmlLength headheight=\the\headheight}\lthtmltypeout{}%
\lthtmltypeout{latex2htmlLength headsep=\the\headsep}\lthtmltypeout{}%
\lthtmltypeout{latex2htmlLength parskip=\the\parskip}\lthtmltypeout{}%
\lthtmltypeout{latex2htmlLength oddsidemargin=\the\oddsidemargin}\lthtmltypeout{}%
\makeatletter
\if@twoside\lthtmltypeout{latex2htmlLength evensidemargin=\the\evensidemargin}%
\else\lthtmltypeout{latex2htmlLength evensidemargin=\the\oddsidemargin}\fi%
\lthtmltypeout{}%
\makeatother
\setcounter{page}{1}
\onecolumn

% !!! IMAGES START HERE !!!

{\newpage\clearpage
\lthtmlinlinemathA{tex2html_wrap_inline1477}%
$ \,\begin{tabular}{|l|}
0  \!\!\! \\
\hline
\end{tabular}\: = \{ \ldots, -15, -10, -5, 0, 5, 10, 15,\ldots \}$%
\lthtmlinlinemathZ
\lthtmlcheckvsize\clearpage}

{\newpage\clearpage
\lthtmlinlinemathA{tex2html_wrap_inline1479}%
$ \,\begin{tabular}{|l|}
1  \!\!\! \\
\hline
\end{tabular}\: = \{ \ldots, -14, -9, -4, 0, 6, 11, 16,\ldots \}$%
\lthtmlinlinemathZ
\lthtmlcheckvsize\clearpage}

{\newpage\clearpage
\lthtmlinlinemathA{tex2html_wrap_inline1481}%
$ \,\begin{tabular}{|l|}
2  \!\!\! \\
\hline
\end{tabular}\: = \{ \ldots, -13, -8, -3, 0, 7, 12, 17,\ldots \}$%
\lthtmlinlinemathZ
\lthtmlcheckvsize\clearpage}

{\newpage\clearpage
\lthtmlinlinemathA{tex2html_wrap_inline1483}%
$ \,\begin{tabular}{|l|}
3  \!\!\! \\
\hline
\end{tabular}\: = \{ \ldots, -12, -7, -2, 0, 8, 13, 18,\ldots \}$%
\lthtmlinlinemathZ
\lthtmlcheckvsize\clearpage}

{\newpage\clearpage
\lthtmlinlinemathA{tex2html_wrap_inline1485}%
$ \,\begin{tabular}{|l|}
4  \!\!\! \\
\hline
\end{tabular}\: = \{ \ldots, -11, -6, -1, 0, 9, 14, 19,\ldots \}$%
\lthtmlinlinemathZ
\lthtmlcheckvsize\clearpage}

{\newpage\clearpage
\lthtmlinlinemathA{tex2html_wrap_inline1487}%
$ a - (a \textrm{ MOD } m)$%
\lthtmlinlinemathZ
\lthtmlcheckvsize\clearpage}

{\newpage\clearpage
\lthtmlinlinemathA{tex2html_wrap_inline1489}%
$ m$%
\lthtmlinlinemathZ
\lthtmlcheckvsize\clearpage}

{\newpage\clearpage
\lthtmlinlinemathA{tex2html_wrap_inline1491}%
$ q \in \mathbb{Z}$%
\lthtmlinlinemathZ
\lthtmlcheckvsize\clearpage}

{\newpage\clearpage
\lthtmlinlinemathA{tex2html_wrap_inline1493}%
$ a = qm + (a \textrm{ MOD } m)$%
\lthtmlinlinemathZ
\lthtmlcheckvsize\clearpage}

{\newpage\clearpage
\lthtmlinlinemathA{tex2html_wrap_inline1495}%
$ a - (a \textrm{ MOD } m) = qm$%
\lthtmlinlinemathZ
\lthtmlcheckvsize\clearpage}

{\newpage\clearpage
\lthtmlinlinemathA{tex2html_wrap_inline1499}%
$ a \textrm{ MOD } m \equiv a ({\rm mod}\  m)$%
\lthtmlinlinemathZ
\lthtmlcheckvsize\clearpage}

{\newpage\clearpage
\lthtmlinlinemathA{tex2html_wrap_inline1501}%
$ \mathbb{Z}_m = \{ \,\begin{tabular}{|l|}
a  \!\!\! \\
\hline
\end{tabular}\: \mid a \in \mathbb{Z} \}$%
\lthtmlinlinemathZ
\lthtmlcheckvsize\clearpage}

{\newpage\clearpage
\lthtmlinlinemathA{tex2html_wrap_inline1503}%
$ \{ \,\begin{tabular}{|l|}
0  \!\!\! \\
\hline
\end{tabular}\:, \,\begin{tabular}{|l|}
1  \!\!\! \\
\hline
\end{tabular}\:, \ldots, \,\begin{tabular}{|l|}
m-1  \!\!\! \\
\hline
\end{tabular}\: \} \subseteq \mathbb{Z}_m$%
\lthtmlinlinemathZ
\lthtmlcheckvsize\clearpage}

{\newpage\clearpage
\lthtmlinlinemathA{tex2html_wrap_inline1505}%
$ a \geqslant m$%
\lthtmlinlinemathZ
\lthtmlcheckvsize\clearpage}

{\newpage\clearpage
\lthtmlinlinemathA{tex2html_wrap_inline1509}%
$ 0 \leqslant r < m$%
\lthtmlinlinemathZ
\lthtmlcheckvsize\clearpage}

{\newpage\clearpage
\lthtmlinlinemathA{tex2html_wrap_inline1511}%
$ a=mq+r$%
\lthtmlinlinemathZ
\lthtmlcheckvsize\clearpage}

{\newpage\clearpage
\lthtmlinlinemathA{tex2html_wrap_inline1513}%
$ a-r$%
\lthtmlinlinemathZ
\lthtmlcheckvsize\clearpage}

{\newpage\clearpage
\lthtmlinlinemathA{tex2html_wrap_inline1517}%
$ a \equiv r ({\rm mod}\  m)$%
\lthtmlinlinemathZ
\lthtmlcheckvsize\clearpage}

{\newpage\clearpage
\lthtmlinlinemathA{tex2html_wrap_inline1519}%
$ \,\begin{tabular}{|l|}
a  \!\!\! \\
\hline
\end{tabular}\: = \,\begin{tabular}{|l|}
r  \!\!\! \\
\hline
\end{tabular}\:$%
\lthtmlinlinemathZ
\lthtmlcheckvsize\clearpage}

{\newpage\clearpage
\lthtmlinlinemathA{tex2html_wrap_inline1521}%
$ \mathbb{Z}_m = \{ \,\begin{tabular}{|l|}
0  \!\!\! \\
\hline
\end{tabular}\:, \,\begin{tabular}{|l|}
1  \!\!\! \\
\hline
\end{tabular}\:, \ldots, \,\begin{tabular}{|l|}
m-1  \!\!\! \\
\hline
\end{tabular}\: \}$%
\lthtmlinlinemathZ
\lthtmlcheckvsize\clearpage}

{\newpage\clearpage
\lthtmlinlinemathA{tex2html_wrap_inline1523}%
$ a \equiv b\  ({\rm mod}\  m)$%
\lthtmlinlinemathZ
\lthtmlcheckvsize\clearpage}

{\newpage\clearpage
\lthtmlinlinemathA{tex2html_wrap_inline1525}%
$ a = b + mq$%
\lthtmlinlinemathZ
\lthtmlcheckvsize\clearpage}

{\newpage\clearpage
\lthtmlinlinemathA{tex2html_wrap_inline1529}%
$ x \equiv y\  ({\rm mod}\  m)$%
\lthtmlinlinemathZ
\lthtmlcheckvsize\clearpage}

{\newpage\clearpage
\lthtmlinlinemathA{tex2html_wrap_inline1531}%
$ x = y + mq'$%
\lthtmlinlinemathZ
\lthtmlcheckvsize\clearpage}

{\newpage\clearpage
\lthtmlinlinemathA{tex2html_wrap_inline1533}%
$ q' \in \mathbb{Z}$%
\lthtmlinlinemathZ
\lthtmlcheckvsize\clearpage}

{\newpage\clearpage
\lthtmlinlinemathA{tex2html_wrap_inline1535}%
$ ax  = (b+mq)(y+mq') = by + bmq' + ymq + mmqq' = by + m(bq' + yq + mqq')$%
\lthtmlinlinemathZ
\lthtmlcheckvsize\clearpage}

{\newpage\clearpage
\lthtmlinlinemathA{tex2html_wrap_inline1537}%
$ ax-by = m(bq'+yq+mqq')$%
\lthtmlinlinemathZ
\lthtmlcheckvsize\clearpage}

{\newpage\clearpage
\lthtmlinlinemathA{tex2html_wrap_inline1539}%
$ ax \equiv by\  {\rm mod}\  m)$%
\lthtmlinlinemathZ
\lthtmlcheckvsize\clearpage}

{\newpage\clearpage
\lthtmlinlinemathA{tex2html_wrap_inline1541}%
$ \,\begin{tabular}{|l|}
a  \!\!\! \\
\hline
\end{tabular}\: = \,\begin{tabular}{|l|}
b  \!\!\! \\
\hline
\end{tabular}\:$%
\lthtmlinlinemathZ
\lthtmlcheckvsize\clearpage}

{\newpage\clearpage
\lthtmlinlinemathA{tex2html_wrap_inline1543}%
$ \,\begin{tabular}{|l|}
x  \!\!\! \\
\hline
\end{tabular}\: = \,\begin{tabular}{|l|}
y  \!\!\! \\
\hline
\end{tabular}\:$%
\lthtmlinlinemathZ
\lthtmlcheckvsize\clearpage}

{\newpage\clearpage
\lthtmlinlinemathA{tex2html_wrap_inline1545}%
$ \,\begin{tabular}{|l|}
ax  \!\!\! \\
\hline
\end{tabular}\: = \,\begin{tabular}{|l|}
by  \!\!\! \\
\hline
\end{tabular}\:$%
\lthtmlinlinemathZ
\lthtmlcheckvsize\clearpage}

{\newpage\clearpage
\lthtmlinlinemathA{tex2html_wrap_inline1547}%
$ \,\begin{tabular}{|l|}
a  \!\!\! \\
\hline
\end{tabular}\: = \,\begin{tabular}{|l|}
b  \!\!\! \\
\hline
\end{tabular}\:
\Rightarrow a \equiv b\  {\rm mod}\  m$%
\lthtmlinlinemathZ
\lthtmlcheckvsize\clearpage}

{\newpage\clearpage
\lthtmlinlinemathA{tex2html_wrap_inline1549}%
$ \,\begin{tabular}{|l|}
x  \!\!\! \\
\hline
\end{tabular}\: = \,\begin{tabular}{|l|}
y  \!\!\! \\
\hline
\end{tabular}\: \Rightarrow
x \equiv y\  {\rm mod}\  m$%
\lthtmlinlinemathZ
\lthtmlcheckvsize\clearpage}

{\newpage\clearpage
\lthtmlinlinemathA{tex2html_wrap_inline1551}%
$ ax \equiv by\  {\rm mod}\  m$%
\lthtmlinlinemathZ
\lthtmlcheckvsize\clearpage}

{\newpage\clearpage
\lthtmlinlinemathA{tex2html_wrap_inline1555}%
$ a,b,c \in \mathbb{Z}$%
\lthtmlinlinemathZ
\lthtmlcheckvsize\clearpage}

{\newpage\clearpage
\lthtmlinlinemathA{tex2html_wrap_inline1557}%
$ (\,\begin{tabular}{|l|}
a  \!\!\! \\
\hline
\end{tabular}\: \odot \,\begin{tabular}{|l|}
b  \!\!\! \\
\hline
\end{tabular}\:) \odot \,\begin{tabular}{|l|}
c  \!\!\! \\
\hline
\end{tabular}\: = \,\begin{tabular}{|l|}
a  \!\!\! \\
\hline
\end{tabular}\: \odot
(\,\begin{tabular}{|l|}
b  \!\!\! \\
\hline
\end{tabular}\: \odot \,\begin{tabular}{|l|}
c  \!\!\! \\
\hline
\end{tabular}\:)$%
\lthtmlinlinemathZ
\lthtmlcheckvsize\clearpage}

{\newpage\clearpage
\lthtmldisplayA{displaymath1559}%
\begin{displaymath}\begin{array}[t]{rcl} (\,\begin{tabular}{|l|}
\!\!  a  \!\!\! \\
\hline
\end{tabular}\: \odot \,\begin{tabular}{|l|}
\!\!  b  \!\!\! \\
\hline
\end{tabular}\:) \odot
\,\begin{tabular}{|l|}
\!\!  c  \!\!\! \\
\hline
\end{tabular}\: & = &
\,\begin{tabular}{|l|}
\!\!  a \cdot b  \!\!\! \\
\hline
\end{tabular}\: \odot \,\begin{tabular}{|l|}
\!\!  c  \!\!\! \\
\hline
\end{tabular}\: \\\\
& = & \,\begin{tabular}{|l|}
\!\!  (a\cdot b) \cdot c  \!\!\! \\
\hline
\end{tabular}\: \\\\
& = & \,\begin{tabular}{|l|}
\!\!  a \cdot (b \cdot c)  \!\!\! \\
\hline
\end{tabular}\: \\\\
& = & \,\begin{tabular}{|l|}
\!\!  a  \!\!\! \\
\hline
\end{tabular}\: \odot \,\begin{tabular}{|l|}
\!\!  b \cdot c  \!\!\! \\
\hline
\end{tabular}\: \\\\
& = & \,\begin{tabular}{|l|}
\!\!  a  \!\!\! \\
\hline
\end{tabular}\: \odot (\,\begin{tabular}{|l|}
\!\!  b  \!\!\! \\
\hline
\end{tabular}\: \odot \,\begin{tabular}{|l|}
\!\!  c  \!\!\! \\
\hline
\end{tabular}\:).
\end{array}\end{displaymath}%
\lthtmldisplayZ
\lthtmlcheckvsize\clearpage}

{\newpage\clearpage
\lthtmlinlinemathA{tex2html_wrap_inline1561}%
$ \,\begin{tabular}{|l|}
a  \!\!\! \\
\hline
\end{tabular}\: \odot \,\begin{tabular}{|l|}
1  \!\!\! \\
\hline
\end{tabular}\: = \,\begin{tabular}{|l|}
a \cdot 1  \!\!\! \\
\hline
\end{tabular}\: =
\,\begin{tabular}{|l|}
a  \!\!\! \\
\hline
\end{tabular}\:$%
\lthtmlinlinemathZ
\lthtmlcheckvsize\clearpage}

{\newpage\clearpage
\lthtmlinlinemathA{tex2html_wrap_inline1563}%
$ \,\begin{tabular}{|l|}
1  \!\!\! \\
\hline
\end{tabular}\: \odot \,\begin{tabular}{|l|}
a  \!\!\! \\
\hline
\end{tabular}\: = \,\begin{tabular}{|l|}
1 \cdot a  \!\!\! \\
\hline
\end{tabular}\: =
\,\begin{tabular}{|l|}
a  \!\!\! \\
\hline
\end{tabular}\:$%
\lthtmlinlinemathZ
\lthtmlcheckvsize\clearpage}

{\newpage\clearpage
\lthtmlinlinemathA{tex2html_wrap_inline1565}%
$ \,\begin{tabular}{|l|}
1  \!\!\! \\
\hline
\end{tabular}\:$%
\lthtmlinlinemathZ
\lthtmlcheckvsize\clearpage}

{\newpage\clearpage
\lthtmlinlinemathA{tex2html_wrap_inline1567}%
$ \langle \mathbb{Z}_m, \odot \rangle$%
\lthtmlinlinemathZ
\lthtmlcheckvsize\clearpage}

{\newpage\clearpage
\lthtmldisplayA{displaymath1569}%
\begin{displaymath}\begin{array}[t]{rcl}
1191 & = & (1) \cdot 1050 + 141 \\
1050 & = & (7) \cdot  141 +  63 \\
141  & = & (2) \cdot   63 +  15 \\
63   & = & (4) \cdot   15 +   3 \\
15   & = & (5) \cdot    3 +   0
\end{array}\end{displaymath}%
\lthtmldisplayZ
\lthtmlcheckvsize\clearpage}

{\newpage\clearpage
\lthtmlinlinemathA{tex2html_wrap_inline1571}%
$ (1191, 1050) = 1$%
\lthtmlinlinemathZ
\lthtmlcheckvsize\clearpage}

{\newpage\clearpage
\lthtmlinlinemathA{tex2html_wrap_inline1573}%
$ s,t \in \mathbb{Z}$%
\lthtmlinlinemathZ
\lthtmlcheckvsize\clearpage}

{\newpage\clearpage
\lthtmlinlinemathA{tex2html_wrap_inline1575}%
$ 1=1191S + 1050t$%
\lthtmlinlinemathZ
\lthtmlcheckvsize\clearpage}

{\newpage\clearpage
\lthtmldisplayA{displaymath1577}%
\begin{displaymath}\begin{array}[t]{rcl}
3   & = & 63 - (4) \cdot 15                                          \\
    & = & 63 - (4) \cdot [141 - (2) \cdot 63 ]                       \\
    & = & (-4) \cdot 141 + (9) \cdot 63                              \\
    & = & (-4) \cdot 141 + (9) \cdot [1050 - (7) \cdot 141 ]         \\
    & = & (9) \cdot 1050 + (-67) \cdot 141                           \\
    & = & (9) \cdot 1050 + (-67) \cdot [1191 - (1) \cdot 1050 ]      \\
    & = & (-67) \cdot 1191 + (76) \cdot 1050
\end{array}\end{displaymath}%
\lthtmldisplayZ
\lthtmlcheckvsize\clearpage}

{\newpage\clearpage
\lthtmlinlinemathA{tex2html_wrap_inline1579}%
$ s=-67$%
\lthtmlinlinemathZ
\lthtmlcheckvsize\clearpage}

{\newpage\clearpage
\lthtmlinlinemathA{tex2html_wrap_inline1581}%
$ t=76$%
\lthtmlinlinemathZ
\lthtmlcheckvsize\clearpage}

{\newpage\clearpage
\lthtmlinlinemathA{tex2html_wrap_inline1583}%
$ (a,b)=1 \Rightarrow (\exists
s,t)(as+bt=1)$%
\lthtmlinlinemathZ
\lthtmlcheckvsize\clearpage}

{\newpage\clearpage
\lthtmlinlinemathA{tex2html_wrap_inline1585}%
$ d=1$%
\lthtmlinlinemathZ
\lthtmlcheckvsize\clearpage}

{\newpage\clearpage
\lthtmlinlinemathA{tex2html_wrap_inline1587}%
$ s$%
\lthtmlinlinemathZ
\lthtmlcheckvsize\clearpage}

{\newpage\clearpage
\lthtmlinlinemathA{tex2html_wrap_inline1589}%
$ t$%
\lthtmlinlinemathZ
\lthtmlcheckvsize\clearpage}

{\newpage\clearpage
\lthtmlinlinemathA{tex2html_wrap_inline1591}%
$ 1 = as+bt$%
\lthtmlinlinemathZ
\lthtmlcheckvsize\clearpage}

{\newpage\clearpage
\lthtmlinlinemathA{tex2html_wrap_inline1593}%
$ (a,b)=1$%
\lthtmlinlinemathZ
\lthtmlcheckvsize\clearpage}

{\newpage\clearpage
\lthtmlinlinemathA{tex2html_wrap_inline1595}%
$ a$%
\lthtmlinlinemathZ
\lthtmlcheckvsize\clearpage}

{\newpage\clearpage
\lthtmlinlinemathA{tex2html_wrap_inline1597}%
$ b$%
\lthtmlinlinemathZ
\lthtmlcheckvsize\clearpage}

{\newpage\clearpage
\lthtmlinlinemathA{tex2html_wrap_inline1599}%
$ d \in
\mathbb{N}$%
\lthtmlinlinemathZ
\lthtmlcheckvsize\clearpage}

{\newpage\clearpage
\lthtmlinlinemathA{tex2html_wrap_inline1601}%
$ d \vert a$%
\lthtmlinlinemathZ
\lthtmlcheckvsize\clearpage}

{\newpage\clearpage
\lthtmlinlinemathA{tex2html_wrap_inline1603}%
$ d \vert b$%
\lthtmlinlinemathZ
\lthtmlcheckvsize\clearpage}

{\newpage\clearpage
\lthtmlinlinemathA{tex2html_wrap_inline1605}%
$ d \vert
as$%
\lthtmlinlinemathZ
\lthtmlcheckvsize\clearpage}

{\newpage\clearpage
\lthtmlinlinemathA{tex2html_wrap_inline1607}%
$ s \in \mathbb{Z}$%
\lthtmlinlinemathZ
\lthtmlcheckvsize\clearpage}

{\newpage\clearpage
\lthtmlinlinemathA{tex2html_wrap_inline1609}%
$ d \vert bt$%
\lthtmlinlinemathZ
\lthtmlcheckvsize\clearpage}

{\newpage\clearpage
\lthtmlinlinemathA{tex2html_wrap_inline1611}%
$ t \in
\mathbb{Z}$%
\lthtmlinlinemathZ
\lthtmlcheckvsize\clearpage}

{\newpage\clearpage
\lthtmlinlinemathA{tex2html_wrap_inline1613}%
$ d \vert (as+bt)$%
\lthtmlinlinemathZ
\lthtmlcheckvsize\clearpage}

{\newpage\clearpage
\lthtmlinlinemathA{tex2html_wrap_inline1615}%
$ d \vert
1$%
\lthtmlinlinemathZ
\lthtmlcheckvsize\clearpage}

{\newpage\clearpage
\lthtmlinlinemathA{tex2html_wrap_inline1617}%
$ d \leqslant 1$%
\lthtmlinlinemathZ
\lthtmlcheckvsize\clearpage}

{\newpage\clearpage
\lthtmlinlinemathA{tex2html_wrap_inline1626}%
$ \mathbf{(a \Rightarrow b)}$%
\lthtmlinlinemathZ
\lthtmlcheckvsize\clearpage}

{\newpage\clearpage
\lthtmlinlinemathA{tex2html_wrap_inline1628}%
$ (a,m)=1$%
\lthtmlinlinemathZ
\lthtmlcheckvsize\clearpage}

{\newpage\clearpage
\lthtmlinlinemathA{tex2html_wrap_inline1632}%
$ 1=as+mt$%
\lthtmlinlinemathZ
\lthtmlcheckvsize\clearpage}

{\newpage\clearpage
\lthtmlinlinemathA{tex2html_wrap_inline1634}%
$ 1 \equiv as+mt ({\rm mod}\  m)$%
\lthtmlinlinemathZ
\lthtmlcheckvsize\clearpage}

{\newpage\clearpage
\lthtmlinlinemathA{tex2html_wrap_inline1636}%
$ 1 \equiv as ({\rm mod}\  m)$%
\lthtmlinlinemathZ
\lthtmlcheckvsize\clearpage}

{\newpage\clearpage
\lthtmlinlinemathA{tex2html_wrap_inline1638}%
$ m \equiv 0 ({\rm
mod}\  m)$%
\lthtmlinlinemathZ
\lthtmlcheckvsize\clearpage}

{\newpage\clearpage
\lthtmlinlinemathA{tex2html_wrap_inline1640}%
$ \mathbf{(b \Rightarrow c)}$%
\lthtmlinlinemathZ
\lthtmlcheckvsize\clearpage}

{\newpage\clearpage
\lthtmlinlinemathA{tex2html_wrap_inline1646}%
$ \,\begin{tabular}{|l|}
1  \!\!\! \\
\hline
\end{tabular}\: = \,\begin{tabular}{|l|}
as  \!\!\! \\
\hline
\end{tabular}\: = \,\begin{tabular}{|l|}
a  \!\!\! \\
\hline
\end{tabular}\: \odot \,\begin{tabular}{|l|}
s  \!\!\! \\
\hline
\end{tabular}\:$%
\lthtmlinlinemathZ
\lthtmlcheckvsize\clearpage}

{\newpage\clearpage
\lthtmlinlinemathA{tex2html_wrap_inline1648}%
$ \,\begin{tabular}{|l|}
a  \!\!\! \\
\hline
\end{tabular}\:$%
\lthtmlinlinemathZ
\lthtmlcheckvsize\clearpage}

{\newpage\clearpage
\lthtmlinlinemathA{tex2html_wrap_inline1650}%
$ <\mathbb{Z}_m, \odot
>$%
\lthtmlinlinemathZ
\lthtmlcheckvsize\clearpage}

{\newpage\clearpage
\lthtmlinlinemathA{tex2html_wrap_inline1652}%
$ \,\begin{tabular}{|l|}
s  \!\!\! \\
\hline
\end{tabular}\:$%
\lthtmlinlinemathZ
\lthtmlcheckvsize\clearpage}

{\newpage\clearpage
\lthtmlinlinemathA{tex2html_wrap_inline1654}%
$ \mathbf{(c \Rightarrow a)}$%
\lthtmlinlinemathZ
\lthtmlcheckvsize\clearpage}

{\newpage\clearpage
\lthtmlinlinemathA{tex2html_wrap_inline1662}%
$ \,\begin{tabular}{|l|}
a  \!\!\! \\
\hline
\end{tabular}\: \odot \,\begin{tabular}{|l|}
s  \!\!\! \\
\hline
\end{tabular}\: = \,\begin{tabular}{|l|}
1  \!\!\! \\
\hline
\end{tabular}\:$%
\lthtmlinlinemathZ
\lthtmlcheckvsize\clearpage}

{\newpage\clearpage
\lthtmlinlinemathA{tex2html_wrap_inline1664}%
$ \,\begin{tabular}{|l|}
as  \!\!\! \\
\hline
\end{tabular}\: =
\,\begin{tabular}{|l|}
1  \!\!\! \\
\hline
\end{tabular}\:$%
\lthtmlinlinemathZ
\lthtmlcheckvsize\clearpage}

{\newpage\clearpage
\lthtmlinlinemathA{tex2html_wrap_inline1666}%
$ as \equiv 1 ({\rm mod}\  m)$%
\lthtmlinlinemathZ
\lthtmlcheckvsize\clearpage}

{\newpage\clearpage
\lthtmlinlinemathA{tex2html_wrap_inline1670}%
$ mt \equiv 0 ({\rm mod}\  m)$%
\lthtmlinlinemathZ
\lthtmlcheckvsize\clearpage}

{\newpage\clearpage
\lthtmlinlinemathA{tex2html_wrap_inline1674}%
$ as+mt \equiv 1 ({\rm mod}\  m)$%
\lthtmlinlinemathZ
\lthtmlcheckvsize\clearpage}

{\newpage\clearpage
\lthtmlinlinemathA{tex2html_wrap_inline1680}%
$ \mathbb{Z}_{12}$%
\lthtmlinlinemathZ
\lthtmlcheckvsize\clearpage}

{\newpage\clearpage
\lthtmlinlinemathA{tex2html_wrap_inline1682}%
$ (a,12)=1$%
\lthtmlinlinemathZ
\lthtmlcheckvsize\clearpage}

{\newpage\clearpage
\lthtmlinlinemathA{tex2html_wrap_inline1688}%
$ \,\begin{tabular}{|l|}
5  \!\!\! \\
\hline
\end{tabular}\:$%
\lthtmlinlinemathZ
\lthtmlcheckvsize\clearpage}

{\newpage\clearpage
\lthtmlinlinemathA{tex2html_wrap_inline1690}%
$ \,\begin{tabular}{|l|}
7  \!\!\! \\
\hline
\end{tabular}\:$%
\lthtmlinlinemathZ
\lthtmlcheckvsize\clearpage}

{\newpage\clearpage
\lthtmlinlinemathA{tex2html_wrap_inline1692}%
$ \,\begin{tabular}{|l|}
11  \!\!\! \\
\hline
\end{tabular}\:$%
\lthtmlinlinemathZ
\lthtmlcheckvsize\clearpage}

{\newpage\clearpage
\lthtmlinlinemathA{tex2html_wrap_inline1694}%
$ p$%
\lthtmlinlinemathZ
\lthtmlcheckvsize\clearpage}

{\newpage\clearpage
\lthtmlinlinemathA{tex2html_wrap_inline1696}%
$ n \in \mathbb{N}$%
\lthtmlinlinemathZ
\lthtmlcheckvsize\clearpage}

{\newpage\clearpage
\lthtmlinlinemathA{tex2html_wrap_inline1698}%
$ 0 < n < p$%
\lthtmlinlinemathZ
\lthtmlcheckvsize\clearpage}

{\newpage\clearpage
\lthtmlinlinemathA{tex2html_wrap_inline1700}%
$ (n,p)=1$%
\lthtmlinlinemathZ
\lthtmlcheckvsize\clearpage}

{\newpage\clearpage
\lthtmlinlinemathA{tex2html_wrap_inline1702}%
$ p-1$%
\lthtmlinlinemathZ
\lthtmlcheckvsize\clearpage}

{\newpage\clearpage
\lthtmlinlinemathA{tex2html_wrap_inline1704}%
$ \mathbb{Z}_{p}$%
\lthtmlinlinemathZ
\lthtmlcheckvsize\clearpage}

{\newpage\clearpage
\lthtmlinlinemathA{tex2html_wrap_inline1708}%
$ p_1,p_2,\ldots,p_n$%
\lthtmlinlinemathZ
\lthtmlcheckvsize\clearpage}

{\newpage\clearpage
\lthtmlinlinemathA{tex2html_wrap_inline1710}%
$ m=p_1p_2 \cdots p_n +1$%
\lthtmlinlinemathZ
\lthtmlcheckvsize\clearpage}

{\newpage\clearpage
\lthtmlinlinemathA{tex2html_wrap_inline1714}%
$ \frac{m}{p_i}
= p_1p_2\cdots p_{i-1}p_{i+1}\cdots p_n + 1$%
\lthtmlinlinemathZ
\lthtmlcheckvsize\clearpage}

{\newpage\clearpage
\lthtmlinlinemathA{tex2html_wrap_inline1716}%
$ i=1,2,\ldots,n$%
\lthtmlinlinemathZ
\lthtmlcheckvsize\clearpage}

{\newpage\clearpage
\lthtmlinlinemathA{tex2html_wrap_inline1720}%
$ p_i$%
\lthtmlinlinemathZ
\lthtmlcheckvsize\clearpage}

{\newpage\clearpage
\lthtmldisplayA{displaymath1728}%
\begin{displaymath}\begin{array}[t]{rcl}
                \varphi(12) & = & \# \{ a \mid \,\begin{tabular}{|l|}
\!\!  a  \!\!\! \\
\hline
\end{tabular}\: \textrm{ est inversible dans } \mathbb{Z}_{12} \} \\
                            & = & \# \{ a \mid 0 < a < 12 \textrm{ et } (a,12) = 1 \} \\
                            & = & \# \{ 1, 5, 7, 11 \} \\
                            & = & 4
                \end{array}\end{displaymath}%
\lthtmldisplayZ
\lthtmlcheckvsize\clearpage}

{\newpage\clearpage
\lthtmldisplayA{displaymath1730}%
\begin{displaymath}\begin{array}[t]{rcl}
                \varphi(p)  & = & \# \{ a \mid \,\begin{tabular}{|l|}
\!\!  a  \!\!\! \\
\hline
\end{tabular}\: \textrm{ est inversible dans } \mathbb{Z}_{p} \} \\
                            & = & \# \{ a \mid 0 < a < p \textrm{ et } (a,p) = 1 \} \\
                            & = & p-1 \textrm{ (en vertu de l'exercice 3.2.13)}
                \end{array}\end{displaymath}%
\lthtmldisplayZ
\lthtmlcheckvsize\clearpage}

{\newpage\clearpage
\lthtmlinlinemathA{tex2html_wrap_inline1735}%
$ \{ 1\cdot p, 2\cdot p, 3\cdot p, \ldots, (q-2)\cdot p, (q-1)\cdot p \}$%
\lthtmlinlinemathZ
\lthtmlcheckvsize\clearpage}

{\newpage\clearpage
\lthtmlinlinemathA{tex2html_wrap_inline1737}%
$ \{ 1\cdot q, 2\cdot q, 3\cdot q, \ldots, (p-2)\cdot q, (p-1)\cdot q \}$%
\lthtmlinlinemathZ
\lthtmlcheckvsize\clearpage}

{\newpage\clearpage
\lthtmlinlinemathA{tex2html_wrap_inline1739}%
$ \# \{ 1\cdot p, 2\cdot p, 3\cdot p, \ldots, (q-2)\cdot p, (q-1)\cdot p \} +
\# \{ 1\cdot q, 2\cdot q, 3\cdot q, \ldots, (p-2)\cdot q, (p-1)\cdot q \} = (q-1) + (p-1)$%
\lthtmlinlinemathZ
\lthtmlcheckvsize\clearpage}

{\newpage\clearpage
\lthtmldisplayA{displaymath1741}%
\begin{displaymath}\begin{array}[t]{rcl}
                \varphi(pq)  & = & \# \{ a \mid \,\begin{tabular}{|l|}
\!\!  a  \!\!\! \\
\hline
\end{tabular}\: \textrm{ est inversible dans } \mathbb{Z}_{pq} \} \\
                             & = & \# \{ a \mid 0 < a < pq \textrm{ et } (a,pq) = 1 \} \\
                             & = & (pq-1) - \# \{ a \mid 0 < a < pq \textrm{ et } (a,pq) \neq 1 \} \\
                             & = & (pq-1) - \# \{ a \mid 0 < a < pq \textrm{ et $a$\  est un } \\
                             &   & \hspace{3cm} \textrm{ multiple de $p$\  ou de $q$\  } \} \\
                             & = & (pq-1) - (q-1) - (p-1) \\
                             & = & pq - q - p + 1 \\
                             & = & q(p-1) - p + 1 \\
                             & = & q(p-1) - (p-1) \\
                             & = & (q-1)(p-1)
                \end{array}\end{displaymath}%
\lthtmldisplayZ
\lthtmlcheckvsize\clearpage}

{\newpage\clearpage
\lthtmlinlinemathA{tex2html_wrap_inline1743}%
$ m=p$%
\lthtmlinlinemathZ
\lthtmlcheckvsize\clearpage}

{\newpage\clearpage
\lthtmlinlinemathA{tex2html_wrap_inline1747}%
$ \varphi(p) = p-1$%
\lthtmlinlinemathZ
\lthtmlcheckvsize\clearpage}

{\newpage\clearpage
\lthtmlinlinemathA{tex2html_wrap_inline1753}%
$ \varphi(m)$%
\lthtmlinlinemathZ
\lthtmlcheckvsize\clearpage}

{\newpage\clearpage
\lthtmlinlinemathA{tex2html_wrap_inline1757}%
$ n$%
\lthtmlinlinemathZ
\lthtmlcheckvsize\clearpage}

{\newpage\clearpage
\lthtmlinlinemathA{tex2html_wrap_inline1759}%
$ 2$%
\lthtmlinlinemathZ
\lthtmlcheckvsize\clearpage}

{\newpage\clearpage
\lthtmlinlinemathA{tex2html_wrap_inline1761}%
$ \sqrt{n}$%
\lthtmlinlinemathZ
\lthtmlcheckvsize\clearpage}

{\newpage\clearpage
\lthtmlinlinemathA{tex2html_wrap_inline1769}%
$ \sqrt{10^{100}} = (10^{100})^{1/2} =
10^{100/2} = 10^{50}$%
\lthtmlinlinemathZ
\lthtmlcheckvsize\clearpage}

{\newpage\clearpage
\lthtmlinlinemathA{tex2html_wrap_inline1771}%
$ 3,17 \times 10^{33}$%
\lthtmlinlinemathZ
\lthtmlcheckvsize\clearpage}

{\newpage\clearpage
\lthtmlinlinemathA{tex2html_wrap_inline1773}%
$ 2,31
\times 10^{23}$%
\lthtmlinlinemathZ
\lthtmlcheckvsize\clearpage}

{\newpage\clearpage
\lthtmlinlinemathA{tex2html_wrap_inline1775}%
$ \log_{b} x$%
\lthtmlinlinemathZ
\lthtmlcheckvsize\clearpage}

{\newpage\clearpage
\lthtmlinlinemathA{tex2html_wrap_inline1779}%
$ x$%
\lthtmlinlinemathZ
\lthtmlcheckvsize\clearpage}

{\newpage\clearpage
\lthtmlinlinemathA{tex2html_wrap_inline1781}%
$ \log_{10} 1000 = 3$%
\lthtmlinlinemathZ
\lthtmlcheckvsize\clearpage}

{\newpage\clearpage
\lthtmlinlinemathA{tex2html_wrap_inline1783}%
$ \log_{10} 10^{100} =
100$%
\lthtmlinlinemathZ
\lthtmlcheckvsize\clearpage}

{\newpage\clearpage
\lthtmlinlinemathA{tex2html_wrap_inline1785}%
$ \log_{2} 2 = 1$%
\lthtmlinlinemathZ
\lthtmlcheckvsize\clearpage}

{\newpage\clearpage
\lthtmlinlinemathA{tex2html_wrap_inline1787}%
$ \log_{2} 512 = 9$%
\lthtmlinlinemathZ
\lthtmlcheckvsize\clearpage}

{\newpage\clearpage
\lthtmlinlinemathA{tex2html_wrap_inline1789}%
$ \log_{2} 1\;048\;576
= 20$%
\lthtmlinlinemathZ
\lthtmlcheckvsize\clearpage}

{\newpage\clearpage
\lthtmlinlinemathA{tex2html_wrap_inline1791}%
$ \log_{2} 10 \approx 3$%
\lthtmlinlinemathZ
\lthtmlcheckvsize\clearpage}

{\newpage\clearpage
\lthtmlinlinemathA{tex2html_wrap_inline1793}%
$ 2^3 = 8$%
\lthtmlinlinemathZ
\lthtmlcheckvsize\clearpage}

{\newpage\clearpage
\lthtmlinlinemathA{tex2html_wrap_inline1795}%
$ \log_{e} 10
\approx 2$%
\lthtmlinlinemathZ
\lthtmlcheckvsize\clearpage}

{\newpage\clearpage
\lthtmlinlinemathA{tex2html_wrap_inline1797}%
$ e^2 \approx 7,4$%
\lthtmlinlinemathZ
\lthtmlcheckvsize\clearpage}

{\newpage\clearpage
\lthtmlinlinemathA{tex2html_wrap_inline1799}%
$ e^3 \approx 20$%
\lthtmlinlinemathZ
\lthtmlcheckvsize\clearpage}

{\newpage\clearpage
\lthtmlinlinemathA{tex2html_wrap_inline1801}%
$ \log_{b} x^n = n \cdot \log{b} x = nL$%
\lthtmlinlinemathZ
\lthtmlcheckvsize\clearpage}

{\newpage\clearpage
\lthtmlinlinemathA{tex2html_wrap_inline1803}%
$ \log{2} 10^{100} = 100 \log{2} 10 \approx 100 \cdot 3 = 300$%
\lthtmlinlinemathZ
\lthtmlcheckvsize\clearpage}

{\newpage\clearpage
\lthtmlinlinemathA{tex2html_wrap_inline1805}%
$ \ln 10^{100} = 100 \ln 10 \approx 100 \cdot 2 = 200$%
\lthtmlinlinemathZ
\lthtmlcheckvsize\clearpage}

{\newpage\clearpage
\lthtmlinlinemathA{tex2html_wrap_inline1807}%
$ 2 \log_{2}
a$%
\lthtmlinlinemathZ
\lthtmlcheckvsize\clearpage}

{\newpage\clearpage
\lthtmlinlinemathA{tex2html_wrap_inline1809}%
$ a=10^{100}$%
\lthtmlinlinemathZ
\lthtmlcheckvsize\clearpage}

{\newpage\clearpage
\lthtmlinlinemathA{tex2html_wrap_inline1811}%
$ 2 \log_{2} a = 2 \log_{2}
10^{100} = 2 \cdot 100 \cdot \log_{2} 10 = 200 \cdot \log_{2} 10
\approx 200 \cdot 3 = 600$%
\lthtmlinlinemathZ
\lthtmlcheckvsize\clearpage}

{\newpage\clearpage
\lthtmlinlinemathA{tex2html_wrap_inline1813}%
$ a^ma^n=a^{m+n}$%
\lthtmlinlinemathZ
\lthtmlcheckvsize\clearpage}

{\newpage\clearpage
\lthtmlinlinemathA{tex2html_wrap_inline1815}%
$ 27^{2^{j+1}} \textrm{ MOD }
100 = 27^{2^{j}+2^{j}} \textrm{ MOD } 100 = 27^{2^j} \cdot
27^{2^j} \textrm{ MOD } 100 = (27^{2^j} \textrm{ MOD } 100) \cdot
(27^{2^j} \textrm{ MOD } 100)  \textrm{ MOD } 100$%
\lthtmlinlinemathZ
\lthtmlcheckvsize\clearpage}

{\newpage\clearpage
\lthtmldisplayA{displaymath1817}%
\begin{displaymath}\begin{array}[t]{rcl}
\overline{(x+y) \textrm{ MOD } m}   & = & \overline{x+y} \\
                                    & = & \overline{x} \oplus \overline{y} \\
                                    & = & \overline{x \textrm{ MOD } m} \oplus \overline{y \textrm{ MOD } m} \\
                                    & = & \overline{(x \textrm{ MOD } m) + (y \textrm{ MOD } m)} \\
                                    & = & \overline{[(x \textrm{ MOD } m) + (y \textrm{ MOD } m)] \textrm{ MOD } m}
\end{array}\end{displaymath}%
\lthtmldisplayZ
\lthtmlcheckvsize\clearpage}

{\newpage\clearpage
\lthtmlinlinemathA{tex2html_wrap_inline1819}%
$ (x+y) \textrm{ MOD } m$%
\lthtmlinlinemathZ
\lthtmlcheckvsize\clearpage}

{\newpage\clearpage
\lthtmlinlinemathA{tex2html_wrap_inline1821}%
$ [(x \textrm{ MOD } m) + (y
\textrm{ MOD } m)] \textrm{ MOD } m$%
\lthtmlinlinemathZ
\lthtmlcheckvsize\clearpage}

{\newpage\clearpage
\lthtmlinlinemathA{tex2html_wrap_inline1825}%
$ \{0,1,2, \ldots, m-1 \}$%
\lthtmlinlinemathZ
\lthtmlcheckvsize\clearpage}

{\newpage\clearpage
\lthtmldisplayA{displaymath1827}%
\begin{displaymath}\begin{array}[t]{rcl}
                \varphi(63) & = & \# \{ a \mid \,\begin{tabular}{|l|}
\!\!  a  \!\!\! \\
\hline
\end{tabular}\: \textrm{ est inversible dans } \mathbb{Z}_{63} \} \\
                            & = & \# \{ a \mid 0 < a < 63 \textrm{ et } (a,63) = 1 \} \\
                            & = & \# \{ 1, 2, 4, 5, 8, 10, 11, 13, 16, 17, 19, 20, 22, 23, 25, 26, 29, 31, 32, 34, \\
                            &   & \quad  37, 38, 40, 41, 43, 44, 46, 47, 50, 52, 53, 55, 58, 59, 61, 62 \} \\
                            & = & 36
                \end{array}\end{displaymath}%
\lthtmldisplayZ
\lthtmlcheckvsize\clearpage}

{\newpage\clearpage
\lthtmlinlinemathA{tex2html_wrap_inline1829}%
$ 36 = 32 + 4 = 2^5 + 2^2$%
\lthtmlinlinemathZ
\lthtmlcheckvsize\clearpage}

{\newpage\clearpage
\lthtmldisplayA{displaymath1831}%
\begin{displaymath}\begin{array}[t]{c|ccccc}
i & 44^{2^i} & & & & 44^{2^i} \textrm{ MOD } 100 \\
\hline
0 & 44^{1}  & \textrm{MOD} & 63 & = & 44 \\
1 & 44^{2}  & \textrm{MOD} & 63 & = & 46 \\
2 & 46^{2}  & \textrm{MOD} & 63 & = & 37 \\
3 & 37^{2}  & \textrm{MOD} & 63 & = & 46 \\
4 & 46^{2}  & \textrm{MOD} & 63 & = & 37 \\
5 & 37^{2}  & \textrm{MOD} & 63 & = & 46
\end{array}\end{displaymath}%
\lthtmldisplayZ
\lthtmlcheckvsize\clearpage}

{\newpage\clearpage
\lthtmldisplayA{displaymath1833}%
\begin{displaymath}\begin{array}[t]{rcl}
        44^{\varphi(63)} \textrm{ MOD } 63  & = & 44^{\varphi(63)} \textrm{ MOD } 63 \\
                                            & = & 44^{36} \textrm{ MOD } 63 \\
                                            & = & 44^{2^5 + 2^2} \textrm{ MOD } 63 \\
                                            & = & 44^{2^5} \cdot 44^{2^2} \textrm{ MOD } 63 \\
                                            & = & [(44^{2^5} \textrm{ MOD } 63) \cdot (44^{2^2} \textrm{ MOD } 63)] \textrm{ MOD } 63 \\
                                            & = & [46 \cdot 37] \textrm{ MOD } 63 \\
                                            & = & 1702 \textrm{ MOD } 63 \\
                                            & = & 27 \cdot 63 + 1 \textrm{ MOD } 63 \\
                                            & = & 1.
        \end{array}\end{displaymath}%
\lthtmldisplayZ
\lthtmlcheckvsize\clearpage}

{\newpage\clearpage
\lthtmlinlinemathA{tex2html_wrap_inline1835}%
$ 44^{\varphi(63)} \textrm{ MOD } 63 = 1$%
\lthtmlinlinemathZ
\lthtmlcheckvsize\clearpage}

{\newpage\clearpage
\lthtmlinlinemathA{tex2html_wrap_inline1837}%
$ (44,63)=1$%
\lthtmlinlinemathZ
\lthtmlcheckvsize\clearpage}

{\newpage\clearpage
\lthtmlinlinemathA{tex2html_wrap_inline1839}%
$ 10^{100}$%
\lthtmlinlinemathZ
\lthtmlcheckvsize\clearpage}

{\newpage\clearpage
\lthtmlinlinemathA{tex2html_wrap_inline1841}%
$ \log_{e} 10^{100} = 100 \cdot \log_{e} 10 \approx 100 \cdot 2 =
200$%
\lthtmlinlinemathZ
\lthtmlcheckvsize\clearpage}

{\newpage\clearpage
\lthtmlinlinemathA{tex2html_wrap_inline1843}%
$ 10^{101}$%
\lthtmlinlinemathZ
\lthtmlcheckvsize\clearpage}

{\newpage\clearpage
\lthtmlinlinemathA{tex2html_wrap_inline1845}%
$ \log_{e} 10^{101} = 101 \cdot
\log_{e} 10 \approx 101 \cdot 2 = 202$%
\lthtmlinlinemathZ
\lthtmlcheckvsize\clearpage}

{\newpage\clearpage
\lthtmlinlinemathA{tex2html_wrap_inline1851}%
$ \pi(10^{101}) - \pi(10^{100})$%
\lthtmlinlinemathZ
\lthtmlcheckvsize\clearpage}

{\newpage\clearpage
\lthtmldisplayA{displaymath1853}%
\begin{displaymath}\begin{array}[t]{rcl}
\pi(10^{101}) - \pi(10^{100})
    & \approx   & \frac{10^{101}}{\log_{e} 10^{101}} - \frac{10^{100}}{\log_{e} 10^{100}} \\\\
    & =         & \frac{10^{101} \cdot \log_{e} 10^{100} - 10^{100} \cdot \log_{e} 10^{101}}{\log_{e}{10^{100}} \cdot \log_{e}{10^{101}}} \\\\
    & =         & \frac{10^{101} \cdot 200 - 10^{100} \cdot 202}{200 \cdot 202} \\\\
    & =         & \frac{10 \cdot 10^{100} \cdot 200 -  10^{100} \cdot 202}{200 \cdot 202} \\\\
    & =         & \frac{10^{100} \cdot 2000 - 10^{100} \cdot 202}{200 \cdot 202} \\\\
    & =         & \frac{10^{100} (2000 - 202)}{200 \cdot 202} \\\\
    & =         & \frac{10^{100} (1798)}{200 \cdot 202} \\\\
    & \approx   & 10^{100} \cdot 4 \times 10^{-2} \\\\
    & \approx   & 4 \times 10^{98}
\end{array}\end{displaymath}%
\lthtmldisplayZ
\lthtmlcheckvsize\clearpage}

{\newpage\clearpage
\lthtmlinlinemathA{tex2html_wrap_inline1855}%
$ 10^{101}-10^{100}$%
\lthtmlinlinemathZ
\lthtmlcheckvsize\clearpage}

{\newpage\clearpage
\lthtmldisplayA{displaymath1863}%
\begin{displaymath}\begin{array}[t]{rcl}
\frac{10^{101}-10^{100}}{4 \times 10^{98}}  & = & \frac{10 \cdot 10^{100}-10^{100}}{4 \times 10^{98}} \\\\
                                            & = & \frac{10^{100}(10-1)}{4 \times 10^{98}} \\\\
                                            & = & \frac{9 \times 10^{100}}{4 \times 10^{98}} \\\\
                                            & = & \frac{9}{4} \times \frac{10^{100}}{10^{98}} \\\\
                                            & = & 2,25 \times 10^{2} \\\\
                                            & = & 225
\end{array}\end{displaymath}%
\lthtmldisplayZ
\lthtmlcheckvsize\clearpage}

{\newpage\clearpage
\lthtmlinlinemathA{tex2html_wrap_inline1865}%
$ \frac{225}{2}
\approx 110$%
\lthtmlinlinemathZ
\lthtmlcheckvsize\clearpage}

{\newpage\clearpage
\lthtmlinlinemathA{tex2html_wrap_inline1869}%
$ 20(\log_{2} p)^3$%
\lthtmlinlinemathZ
\lthtmlcheckvsize\clearpage}

{\newpage\clearpage
\lthtmlinlinemathA{tex2html_wrap_inline1875}%
$ 20(\log_{2} p)^3 =
20(\log_{2} 10^{100})^3 = 20 (100 \log_{2} 10)^3 = 20 ( 100 \cdot
3)^3 = 20 (300)^3 = 20 \cdot 27\;000\;000 = 540\;000\;000$%
\lthtmlinlinemathZ
\lthtmlcheckvsize\clearpage}

{\newpage\clearpage
\lthtmlinlinemathA{tex2html_wrap_inline1900}%
$ 2\log_{2}a = 2 \log_2 10^{100} = 2 \cdot 100 \log_{2} 10 = 200 \log_{2} 10
\approx 200 \cdot 3 = 600$%
\lthtmlinlinemathZ
\lthtmlcheckvsize\clearpage}

{\newpage\clearpage
\lthtmlinlinemathA{tex2html_wrap_inline1906}%
$ as+bt=d$%
\lthtmlinlinemathZ
\lthtmlcheckvsize\clearpage}

{\newpage\clearpage
\lthtmlinlinemathA{tex2html_wrap_inline1912}%
$ 8\log_{2} p$%
\lthtmlinlinemathZ
\lthtmlcheckvsize\clearpage}

{\newpage\clearpage
\lthtmlinlinemathA{tex2html_wrap_inline1914}%
$ a(s+b) + b(t-a) = as+ab+bt-ab =
as+bt=c$%
\lthtmlinlinemathZ
\lthtmlcheckvsize\clearpage}

{\newpage\clearpage
\lthtmlinlinemathA{tex2html_wrap_inline1916}%
$ a(s+2b)+b(t-2a) = as + 2ab + bt -2ab = as+bt =
c$%
\lthtmlinlinemathZ
\lthtmlcheckvsize\clearpage}

{\newpage\clearpage
\lthtmlinlinemathA{tex2html_wrap_inline1918}%
$ a(s+3b)+b(t-3a) = as + 3ab + bt -3ab = as+bt =
c$%
\lthtmlinlinemathZ
\lthtmlcheckvsize\clearpage}

{\newpage\clearpage
\lthtmlinlinemathA{tex2html_wrap_inline1920}%
$ c$%
\lthtmlinlinemathZ
\lthtmlcheckvsize\clearpage}

{\newpage\clearpage
\lthtmlinlinemathA{tex2html_wrap_inline1922}%
$ k = \Big\lceil  \frac{-s}{b}
\Big\rceil$%
\lthtmlinlinemathZ
\lthtmlcheckvsize\clearpage}

{\newpage\clearpage
\lthtmlinlinemathA{tex2html_wrap_inline1924}%
$ \lceil x \rceil$%
\lthtmlinlinemathZ
\lthtmlcheckvsize\clearpage}

{\newpage\clearpage
\lthtmlinlinemathA{tex2html_wrap_inline1930}%
$ s+\Big\lceil  \frac{-s}{b}
\Big\rceil b$%
\lthtmlinlinemathZ
\lthtmlcheckvsize\clearpage}

{\newpage\clearpage
\lthtmlinlinemathA{tex2html_wrap_inline1934}%
$ t- \Big\lceil  \frac{-s}{b}
\Big\rceil a$%
\lthtmlinlinemathZ
\lthtmlcheckvsize\clearpage}


\end{document}
